\documentclass[11pt, a4paper]{article}\usepackage[]{graphicx}\usepackage[]{xcolor}
% maxwidth is the original width if it is less than linewidth
% otherwise use linewidth (to make sure the graphics do not exceed the margin)
\makeatletter
\def\maxwidth{ %
  \ifdim\Gin@nat@width>\linewidth
    \linewidth
  \else
    \Gin@nat@width
  \fi
}
\makeatother

\definecolor{fgcolor}{rgb}{0.345, 0.345, 0.345}
\newcommand{\hlnum}[1]{\textcolor[rgb]{0.686,0.059,0.569}{#1}}%
\newcommand{\hlsng}[1]{\textcolor[rgb]{0.192,0.494,0.8}{#1}}%
\newcommand{\hlcom}[1]{\textcolor[rgb]{0.678,0.584,0.686}{\textit{#1}}}%
\newcommand{\hlopt}[1]{\textcolor[rgb]{0,0,0}{#1}}%
\newcommand{\hldef}[1]{\textcolor[rgb]{0.345,0.345,0.345}{#1}}%
\newcommand{\hlkwa}[1]{\textcolor[rgb]{0.161,0.373,0.58}{\textbf{#1}}}%
\newcommand{\hlkwb}[1]{\textcolor[rgb]{0.69,0.353,0.396}{#1}}%
\newcommand{\hlkwc}[1]{\textcolor[rgb]{0.333,0.667,0.333}{#1}}%
\newcommand{\hlkwd}[1]{\textcolor[rgb]{0.737,0.353,0.396}{\textbf{#1}}}%
\let\hlipl\hlkwb

\usepackage{framed}
\makeatletter
\newenvironment{kframe}{%
 \def\at@end@of@kframe{}%
 \ifinner\ifhmode%
  \def\at@end@of@kframe{\end{minipage}}%
  \begin{minipage}{\columnwidth}%
 \fi\fi%
 \def\FrameCommand##1{\hskip\@totalleftmargin \hskip-\fboxsep
 \colorbox{shadecolor}{##1}\hskip-\fboxsep
     % There is no \\@totalrightmargin, so:
     \hskip-\linewidth \hskip-\@totalleftmargin \hskip\columnwidth}%
 \MakeFramed {\advance\hsize-\width
   \@totalleftmargin\z@ \linewidth\hsize
   \@setminipage}}%
 {\par\unskip\endMakeFramed%
 \at@end@of@kframe}
\makeatother

\definecolor{shadecolor}{rgb}{.97, .97, .97}
\definecolor{messagecolor}{rgb}{0, 0, 0}
\definecolor{warningcolor}{rgb}{1, 0, 1}
\definecolor{errorcolor}{rgb}{1, 0, 0}
\newenvironment{knitrout}{}{} % an empty environment to be redefined in TeX

\usepackage{alltt}

\usepackage[top = 0.6 in, bottom = 0.6 in, left = 0.8 in, right = 0.8 in]{geometry}
\usepackage{amsmath, amssymb, amsfonts}

\allowdisplaybreaks[4]

\usepackage{enumerate}
\usepackage{array}
\usepackage{multirow}
\usepackage{dingbat}
\usepackage{fontawesome5}
\usepackage{tasks}
\usepackage{bbding}
\usepackage{twemojis}
% how to use bull's eye ----- \scalebox{2.0}{\twemoji{bullseye}}
\usepackage{fontspec}
\usepackage{fancyhdr}
\usepackage{lastpage}

\pagestyle{fancy}
\fancyhf{}
\fancyfoot[C]{Page \thepage\ of \pageref{LastPage}}
\renewcommand{\headrulewidth}{0pt}

\title{MSMS 408 : Practical 03}
\author{{\fontsize{22}{20} \benyorithafont Ananda Biswas} \\[1em] Exam Roll No. : 24419STC053}
\date{\today}

\newfontface\benyorithafont{Benyoritha-G334D.otf}

\newfontface\myfont{Myfont1-Regular.ttf}[LetterSpace=0.05em]

\newfontface\cbfont{CaveatBrush-Regular.ttf}
% how to use --- \myfont --write text here--
\IfFileExists{upquote.sty}{\usepackage{upquote}}{}
\begin{document}

\maketitle


\section*{\faArrowAltCircleRight[regular] \textcolor{blue}{Question}}

\hspace{1cm} Consider a random sample $\{0.6, 1.2, 0.9, 1.8, 0.5, 1.1, 0.7, 1.6, 0.8, 1.3, 0.4, 1.0, 0.9, 1.5, 0.6\}$ drawn from an Exponential distribution with parameter $\lambda$. Assume that the prior distribution of $\lambda$ follows a Gamma distribution with mean $1.5$ and variance $0.75$.

\begin{enumerate}[(i)]
\item Write the prior density and the likelihood function and derive the posterior distribution of $\lambda$.

\item Compute the posterior mean and posterior variance of $\lambda$.

\item Obtain the Bayes estimate of $\lambda$ under the squared error loss function and absolute error loss function.

\item Generate a sample of size $20000$ from the posterior distribution and plot the prior and posterior densities along with a histogram of the generated posterior sample.

\item Provide a $95\%$ Bayesian credible interval for $\lambda$ based on simulated posterior values.

\item Perform a posterior predictive simulation to estimate the probability that a future observation exceeds $2$.
\end{enumerate}

\section*{\faArrowAltCircleRight[regular] \textcolor{blue}{Theory}}

\begin{itemize}

\item[\scalebox{2.0}{\twemoji{keycap: 1}}] \hspace{0.1cm} Let the gamma distribution has shape $ = \alpha$ and rate $= \beta$. Then, 

\begin{center}
mean $ = \dfrac{\alpha}{\beta}$ and variance $ = \dfrac{\alpha}{\beta^2}$.
\end{center}

So $\dfrac{\text{mean}}{\text{variance}} = \beta \Rightarrow \beta = \dfrac{1.5}{0.75} = 2$, consequently $\alpha = \beta \times 1.5 = 3$.\\[0.3em]

So, $\lambda \sim \text{Gamma}(3, 2)$.

The prior distribution is 

$$\lambda \sim g(\lambda) \propto e^{-\beta \lambda} {\lambda}^{\alpha - 1}.$$

The likelihood of $\lambda$ is

$$L(\lambda) \propto {\lambda^n} {e^{-\lambda \sum \limits_{i = 1}^{n} x_i}}.$$

\newpage

$\therefore$ The posterior distribution of $\lambda$ will be

$$g(\lambda | \mathbf{x}) \propto e^{-\beta \lambda - \lambda \sum \limits_{i = 1}^{n} x_i} {\lambda}^{\alpha + n - 1}.$$

Recall that, Gamma is a conjugate prior for Exponential parameter. So we identify the posterior distribution of $\lambda$ as

$$\lambda | \mathbf{x} \sim \text{Gamma}\left( \alpha + n, \beta + \sum \limits_{i = 1}^{n} x_i \right).$$

\begin{knitrout}
\definecolor{shadecolor}{rgb}{0.969, 0.969, 0.969}\color{fgcolor}\begin{kframe}
\begin{alltt}
\hldef{alpha} \hlkwb{<-} \hlnum{3}\hldef{; beta} \hlkwb{<-} \hlnum{2}
\end{alltt}
\end{kframe}
\end{knitrout}

\begin{knitrout}
\definecolor{shadecolor}{rgb}{0.969, 0.969, 0.969}\color{fgcolor}\begin{kframe}
\begin{alltt}
\hldef{x} \hlkwb{<-} \hlkwd{c}\hldef{(}\hlnum{0.6}\hldef{,} \hlnum{1.2}\hldef{,} \hlnum{0.9}\hldef{,} \hlnum{1.8}\hldef{,} \hlnum{0.5}\hldef{,} \hlnum{1.1}\hldef{,} \hlnum{0.7}\hldef{,} \hlnum{1.6}\hldef{,} \hlnum{0.8}\hldef{,} \hlnum{1.3}\hldef{,} \hlnum{0.4}\hldef{,} \hlnum{1.0}\hldef{,} \hlnum{0.9}\hldef{,} \hlnum{1.5}\hldef{,} \hlnum{0.6}\hldef{)}
\end{alltt}
\end{kframe}
\end{knitrout}

From the given, $\lambda | \mathbf{x} \sim \text{Gamma}(18, 16.9)$.

\item[\scalebox{2.0}{\twemoji{keycap: 2}}] \hspace{0.1cm} The posterior mean and variance are

$$E[\lambda | \mathbf{x}] = \dfrac{\alpha + n}{\beta + \sum \limits_{i = 1}^{n} x_i} = 1.065, \qquad \text{Var}[\lambda | \mathbf{x}] = \dfrac{\alpha + n}{ \left( \beta + \sum \limits_{i = 1}^{n} x_i \right)^2} = 0.063.$$

\item[\scalebox{2.0}{\twemoji{keycap: 3}}] \hspace{0.1cm} Under squared error loss function, Bayes estimate of $\lambda$ is the posterior mean $1.065$.\\[0.3em]

Under absolute error loss function, Bayes estimate of $\lambda$ is the posterior median $1.05$.

\end{itemize}

\section*{\faArrowAltCircleRight[regular] \textcolor{blue}{R Program}}





\begin{itemize}
\item[\scalebox{2.0}{\twemoji{keycap: 4}}]
\end{itemize}

\begin{knitrout}
\definecolor{shadecolor}{rgb}{0.969, 0.969, 0.969}\color{fgcolor}\begin{kframe}
\begin{alltt}
\hldef{n} \hlkwb{<-} \hlnum{2} \hlopt{*} \hlnum{10}\hlopt{^}\hlnum{5}
\end{alltt}
\end{kframe}
\end{knitrout}

\begin{knitrout}
\definecolor{shadecolor}{rgb}{0.969, 0.969, 0.969}\color{fgcolor}\begin{kframe}
\begin{alltt}
\hldef{lambda_posterior} \hlkwb{<-} \hlkwd{rgamma}\hldef{(n,} \hlkwc{shape} \hldef{=} \hlnum{18}\hldef{,} \hlkwc{rate} \hldef{=} \hlnum{16.9}\hldef{)}
\end{alltt}
\end{kframe}
\end{knitrout}

\begin{knitrout}
\definecolor{shadecolor}{rgb}{0.969, 0.969, 0.969}\color{fgcolor}\begin{kframe}
\begin{alltt}
\hlkwd{data.frame}\hldef{(lambda_posterior)} \hlopt{%>%}
  \hlkwd{ggplot}\hldef{(}\hlkwd{aes}\hldef{(}\hlkwc{x} \hldef{= lambda_posterior))} \hlopt{+}
  \hlkwd{geom_histogram}\hldef{(}\hlkwd{aes}\hldef{(}\hlkwc{y} \hldef{=} \hlkwd{after_stat}\hldef{(density)),} \hlkwc{bins} \hldef{=} \hlnum{20}\hldef{,}
                 \hlkwc{fill} \hldef{=} \hlsng{"#F54927"}\hldef{,} \hlkwc{color} \hldef{=} \hlsng{"black"}\hldef{,} \hlkwc{linewidth} \hldef{=} \hlnum{1}\hldef{)} \hlopt{+}
  \hlkwd{stat_function}\hldef{(}\hlkwc{fun} \hldef{= dgamma,} \hlkwc{args} \hldef{=} \hlkwd{list}\hldef{(}\hlkwc{shape} \hldef{=} \hlnum{3}\hldef{,} \hlkwc{rate} \hldef{=} \hlnum{2}\hldef{),} \hlkwc{xlim} \hldef{=} \hlkwd{c}\hldef{(}\hlnum{0}\hldef{,} \hlnum{3.5}\hldef{),}
                \hlkwc{linewidth} \hldef{=} \hlnum{1.25}\hldef{,} \hlkwc{linetype} \hldef{=} \hlsng{"dashed"}\hldef{)} \hlopt{+}
  \hlkwd{stat_function}\hldef{(}\hlkwc{fun} \hldef{= dgamma,} \hlkwc{args} \hldef{=} \hlkwd{list}\hldef{(}\hlkwc{shape} \hldef{=} \hlnum{18}\hldef{,} \hlkwc{rate} \hldef{=} \hlnum{16.9}\hldef{),}
                \hlkwc{linewidth} \hldef{=} \hlnum{1.25}\hldef{)} \hlopt{+}
  \hlkwd{labs}\hldef{(}\hlkwc{x} \hldef{=} \hlkwd{expression}\hldef{(lambda),}
       \hlkwc{title} \hldef{=} \hlsng{"Histogram of Simulated Sample from Posterior Distribution"}\hldef{)}
\end{alltt}
\end{kframe}
\includegraphics[width=\maxwidth]{figure/unnamed-chunk-7-1} 
\end{knitrout}


\begin{itemize}
\item[\scalebox{2.0}{\twemoji{keycap: 5}}] $95\%$ Bayesian credible interval based on simulated posterior sample is given by $$(q_{0.025}, q_{0.975})$$ where $q_{\alpha}$ is the $\alpha-$fractile of the simulated sample.

\begin{knitrout}
\definecolor{shadecolor}{rgb}{0.969, 0.969, 0.969}\color{fgcolor}\begin{kframe}
\begin{alltt}
\hlkwd{quantile}\hldef{(lambda_posterior,} \hlkwd{c}\hldef{(}\hlnum{0.025}\hldef{,} \hlnum{0.975}\hldef{))}
\end{alltt}
\begin{verbatim}
##      2.5%     97.5% 
## 0.6326549 1.6101341
\end{verbatim}
\end{kframe}
\end{knitrout}
\end{itemize}

\newpage

\begin{itemize}
\item[\scalebox{2.0}{\twemoji{keycap: 5}}] \hspace{0.3cm} We have to find $$P(X_{\text{new}} > 2 \mid \text{data}).$$

\begin{knitrout}
\definecolor{shadecolor}{rgb}{0.969, 0.969, 0.969}\color{fgcolor}\begin{kframe}
\begin{alltt}
\hldef{l} \hlkwb{<-} \hlkwd{rgamma}\hldef{(n,} \hlkwc{shape} \hldef{=} \hlnum{18}\hldef{,} \hlkwc{rate} \hldef{=} \hlnum{16.9}\hldef{)}

\hldef{x_new} \hlkwb{<-} \hlkwd{rexp}\hldef{(n,} \hlkwc{rate} \hldef{= l)}

\hlkwd{mean}\hldef{(x_new} \hlopt{>} \hlnum{2}\hldef{)}
\end{alltt}
\begin{verbatim}
## [1] 0.13439
\end{verbatim}
\end{kframe}
\end{knitrout}

From simulation, the posterior predictive probability that a future observation exceeds $2$ is $0.13$.

\end{itemize}
\end{document}
